\section{Discussion}\label{sec:discussion}

\paragraph{Why a control plane rather than a better memory.} Under I1, a
sharper memory is also a sharper hazard: the more faithful a retrieved statement
is, the more tempting it becomes to present it as current fact. Moving the
design effort from retrieval quality to admission and degradation addresses the
failure surface directly, at the cost of leaving retrieval quality roughly where
it was.

\paragraph{Tension 1: a fail-open majority with a single fail-closed locus.}
The asymmetry is deliberate---a safety failure must block, a functional failure
must stay usable---but it has a price: functional failures can be silent. The
compensation is observability rather than symmetry. Every dropped block carries
a reason code, downgrades are recorded rather than raised, and each turn emits
an event that can be replayed offline. This does not make silent failure
impossible; it makes it attributable after the fact.

\paragraph{Tension 2: a frozen instrument versus a learning loop.} If the
evaluation instrument is writable, the first thing a learning loop learns is to
edit the exam. Freezing metric definitions, replay set, and guardrail thresholds
prevents that, but shrinks the space in which learning can help. Our own search
illustrates the price: across 972 parameter configurations, no configuration
improved the objective by more than the practical-significance threshold
($\Delta < 5\times10^{-4}$), while all six level-1 retrieval parameters became
identifiable. We report it as a methodological result about the instrument, not
about performance.

\paragraph{Tension 3: self-contained plugin versus a single source of truth.} A
host plugin must be installable without a Python runtime, so the control
parameters exist in three places: the authoritative rules file, the host patch
configuration, and the plugin defaults. A 66-key comparison detects drift, but
detecting drift is not proving behavioural equivalence. The honest description
of the port is a re-implementation with a documented divergence list, not a
binding.

\paragraph{Costs.} Four costs are visible in the current artefacts. First,
budget: injected blocks compete with conversation history, and the trade is
governed by static weights that are themselves guesses. Second, latency: the
control plane is designed to stay within a small self-imposed budget, but that
budget is a target, not a measurement. Third, maintenance: two implementations
of the same semantics must be kept aligned by hand. Fourth, audit volume: 308
audit events over eight real turns is a nontrivial storage and inspection load.

\paragraph{Transferability.} Two layers should be separated. The mechanism layer
---threshold semantics, failure semantics, the frozen-instrument discipline---is
host-agnostic and can be specified independently of any API. The wiring layer
---section-based prompt injection, hook names, the approval seam---is specific
to one host. Porting therefore requires re-deriving the wiring and re-checking
each documented divergence against the new host.

\paragraph{Reachability versus effect size.} The evidence in this paper
establishes three things: the mechanisms are reachable in a real host, the
contract executes, and failures are auditable and replayable. It does not
establish a significant positive effect size on large-scale real traffic. That
gap is structural rather than an oversight: measuring such an effect requires
traffic we cannot manufacture and interventions we should not run on someone
else's deployment. It is precisely why the artefacts are open---the engine, the
replay set, the chaos corpora, the pre-registered protocols and the analysis
tools---so that the measurement can be reproduced or refuted by readers with
better access. We would rather publish the instrument and the measured zeros
than a number we cannot defend.

\paragraph{Methodological lesson.} Treating ``what we cannot claim'' as a
first-class artefact turned out to be cheap and repeatedly useful: an
assumptions ledger, a provenance tag on every number, and a frozen instrument
made overclaiming a mechanical rather than a rhetorical failure. We recommend
the practice independently of this system.
